% Related work: absorb all 13 mechanism families, retain narrow residual.
ORION Paper III's contribution is deliberately narrow. Rather than claiming
novelty in any single mechanism --- extraction, schema matching, data
integration, contradiction detection, or literature-based discovery --- we
situate each as absorbed prior work and identify the precise residual.

\subsection{Structured scientific knowledge bases and extraction}

MUSE~\cite{muse2026} constructs a full-text cross-domain knowledge base of
scientific problems, solutions, and rationales, grounding each structure in
its originating sentences. Through this mechanism, source-grounded
cross-domain structured knowledge extraction is absorbed as prior art.
However, MUSE produces a single problem/solution/rationale structure per
source; it does not distinguish referent, construct, measurement, context,
modality, and attribution as separate typed hypotheses, nor does it
represent the failure to align as an obstruction. The Discovery
Engine~\cite{discoveryEngine2025} similarly builds structured scientific
knowledge artifacts from literature. Both are valuable inputs and baselines
for ORION, not competitors to be contested.

SciER~\cite{scier2024} provides full-text scientific entity and relation
annotations and establishes the extraction quality bar for datasets,
methods, and tasks in scientific documents. Scientific information
extraction, discourse parsing, coreference resolution, and modality
detection are absorbed mechanisms. Its flat entity/relation extraction does
not, by itself, represent the participation of typed coordinates in a
\textsc{glue}-or-obstruction decision; the extraction is a baseline, and the
\emph{vertical} separation of semantic coordinates is residual.

\subsection{Schema-centric approaches}

SciSchema.org~\cite{scischema2026} provides expert-designed
multidisciplinary scientific-process schemas with explicit parameters,
measurements, and provenance --- absorbing multidisciplinary expert schema
design. Its schemas are static per discipline, whereas ORION treats schemas
as \emph{source-local, composable, and obstruable} projections assembled per
research question (residual D1). Expert-designed schemas do not by
themselves guarantee traceability of the path from source to schema
(residual D4).

SCOPE and SCION~\cite{scopes2026} introduce a benchmark and an auditable
reference pipeline for evidence-linked schema induction and fusion. Schema
induction and conservative fusion are absorbed. SCOPE/SCION fuses schemas to
produce a single integrated schema, whereas ORION's \textsc{glue}-or-obstruction
semantics (residual D3) permit fusion only when all relevant coordinates are
compatible; otherwise an obstruction or plural portrait is the correct
output. SCOPE/SCION does not represent measurement/operationalization
equivalence as a separate coordinate or track non-preserved information.

Executable Schema Contracts~\cite{schemaContracts2026} demonstrate
provenance-aware shared-schema integration and retrieval. Provenance-aware
multi-source knowledge integration and retrieval under schema constraints
are absorbed. However, Executable Schema Contracts assume a shared schema
exists or can be derived; ORION allows the case where no shared schema is
justified, producing an obstruction (residual D3), and it extends
recoverability beyond provenance to include \emph{non-preserved
coordinates} (residual D4).

\subsection{Alignment, integration, and resolution}

Ontology and schema matching~\cite{euzenat2007ontology} pairs---including
LLM-based matching such as LLMatch~\cite{llmatch2025}---find correspondences
between schemas or ontologies. The mechanism of finding correspondences is
absorbed. Alignment methods typically produce a single equivalence or
subsumption relation, whereas ORION's D2 distinguishes referent, construct,
measurement, and context as \emph{separate alignment axes}: two schemas may
align at the construct level yet differ at the measurement level. Alignment
methods do not produce preservation conditions (D4) or obstruction verdicts
(D3).

Data integration, entity resolution, and record linkage
\cite{oh2017unified,liu2022scholar} assume a shared referent and produce a
fused record. ORION's D3 admits cases where the correct output is not a
fused record, and D2 distinguishes same-referent/different-construct cases
that data integration collapses into a single alignment score. Federated and
multimodel knowledge graphs \cite{liu2022scholar} treat each graph as
canonical for its domain, whereas ORION treats each source as a local,
non-canonical projection (D1) and tracks what is lost in any mapping (D4).

\subsection{Claim comparison and contradiction}

Contradiction and stance detection \cite{raghunathan2022stance} typically
produce a binary contradiction/non-contradiction or stance label. ORION
distinguishes contradiction from \emph{context difference}, \emph{modality
difference}, \emph{measurement difference}, and \emph{attribution
difference} --- each of which a standard system might label as
``contradiction.'' Only true contradiction (after all coordinates are
aligned) should trigger a revision; the other differences are informative
obstructions or condition-aligned mappings.

\subsection{Measurement equivalance and construct validity}

The psychometric and metrological literatures on measurement equivalence,
construct validity, and invariance testing \cite{sebastian2017measurement}
have long established that constructs and their operationalizations are
distinct. ORION does not claim to discover this fact; it claims to
represent it as a first-class typed coordinate in an automated integration
pipeline and to make the \textsc{glue}-or-obstruction decision depend on it. The
measurement-equivalence literature does not provide a generic computational
representation that plugs into a knowledge integration system.

\subsection{Literature-based discovery}

Swanson's A-B-C bridging \cite{swanson1986} and later systems find implicit
connections between literatures that do not cite each other. The mechanism
of discovering bridges is absorbed. ORION adds the requirement that the
meaning of the bridging concept B be checked for compatibility across both
halves of the bridge before the composition is accepted: if B changes
meaning (referent, construct, measurement, or context) between A-B and B-C,
the bridge is an obstruction, not a discovery. This is the residual D2/D3
applied to literature-based discovery.

\subsection{Scientific pluralism}

The philosophy of science literature on scientific pluralism
\cite{kellert2006,cartwright1999,chang2012} has argued for decades that
multiple incompatible scientific representations can each be valid and
that forced unification is harmful. ORION does not claim to discover
pluralism; it claims that an automated integration system can
\emph{represent} pluralism as a first-class output --- \textsc{obstruction} or
\textsc{plural\_view} --- rather than being forced to produce a single
unified answer. The computational implementation of a pluralism principle
is the residual.

\subsection{Scientific RAG and cross-domain synthesis}

OpenScholar~\cite{openScholar2024} synthesizes scientific literature with
retrieval-augmented generation; BioSage~\cite{biosage2025} performs
cross-domain scientific reasoning. Scientific RAG and cross-domain
translation are absorbed mechanisms, and serve as strong baselines in our
evaluation. LLM-based structured extraction, in-context integration, and
prompt-based schema alignment \cite{adias2026} are likewise absorbed.
These systems typically invoke a single prompt or schema to extract and
integrate, implicitly flattening source projections into a canonical
representation. ORION's architectural invariants --- typed coordinates,
\textsc{glue}-or-obstruction, recoverability --- are orthogonal to the prompt
engineering question and are testable as additive components on top of any
backbone. Whether these invariants measurably reduce false integration
relative to flat integration is the empirical question our gold study
answers.

\subsection{Summary of the absorbed boundary and residual}

\begin{table}[t]
\centering
\small
\begin{tabular}{p{0.32\textwidth}p{0.28\textwidth}p{0.28\textwidth}}
\toprule
\textbf{Mechanism family} & \textbf{Absorbed (related work)} & \textbf{ORION residual} \\
\midrule
Source-grounded structured KB & MUSE, Discovery Engine & Source-local addressable projections (D1) \\
Expert scientific schemas & SciSchema.org & Composable, obstruable projections (D1) \\
Schema induction/fusion & SCOPE/SCION & Conditional GLUE; obstruction (D3) \\
Provenance integration & Executable Schema Contracts & Non-preserved coordinates (D4) \\
Scientific IE & SciER & Typed coordinates gate integration (D2) \\
Ontology/schema alignment & Ontology matching, LLMatch & Separate alignment axes (D2) \\
Measurement equivalence & Psychometrics, metrology & First-class computational coordinate (D2) \\
Data integration & Record linkage, data fusion & Non-merge outputs; recoverability (D3,D4) \\
Contradiction detection & Stance/contradiction systems & Coordinate-resolved contradiction \\
Literature-based discovery & Swanson lineage & Meaning-stability check on bridges (D2,D3) \\
Scientific pluralism & Philosophy of science & Computational representation (D3) \\
Scientific RAG / cross-domain & OpenScholar, BioSage & Testable additive invariants \\
FM-based KG representation & Discovery Engine, ADIAS & Architectural invariants vs prompt engineering \\
\bottomrule
\end{tabular}
\caption{The absorbed boundary. Each mechanism family is prior art; the
residual is the last column.}
\label{tab:absorbed}
\end{table}

We do not claim novelty in any row of the middle column. The claimed
contribution is that the combination of typed coordinate distinctions,
conditional \textsc{glue} with typed obstruction and plural portraits,
recoverability, and the W effect is not addressed by any single prior
mechanism, and that this combination can be evaluated as a measurable
reduction in false integration. Table~\ref{tab:absorbed} summarizes the
boundary.
\endinput